\section{Reduction and notation}
\label{sec:reduction-and-notation}

We prove Theorem~\ref{thm:main} by contradiction along sequences.
All constants in an asymptotic argument, including exponents unless
explicitly indicated otherwise, are independent of \(n\).

\begin{lemma}[Counterexample sequence]\label{lem:counterexample-sequence}
If \(R(Q_n)/2^n\) is unbounded, there is a sequence of dimensions and
red-blue colorings between disjoint sets \(X,Y\), each of size \(N\),
with no monochromatic \(Q_n\) across the two sides, such that
\[
 n\longrightarrow\infty,\qquad C_n=N/2^n\longrightarrow\infty,\qquad N\le n2^n.
\]
\end{lemma}
\begin{proof}
Each fixed dimension has finite Ramsey number by the finite Ramsey
theorem~\cite{ramsey1930}. Thus an unbounded ratio has a subsequence along
which both the dimensions and the ratios tend to infinity. For each such
\(n\), take a coloring of \(K_{R(Q_n)-1}\) containing no monochromatic
\(Q_n\), and set
\[
N=\min\left\{\left\lfloor\frac{R(Q_n)-1}{2}\right\rfloor,
                 n2^n\right\}.
\]
Choose \(2N\) of its vertices and split them into equal sides. Restriction
to the cross edges cannot create a monochromatic cube. Moreover,
\(N/2^n\to\infty\), since both terms in the defining minimum have ratios
tending to infinity.
\end{proof}

Fix such a sequence for the remainder of the proof. We will contradict its
existence without imposing any rate of growth on \(C_n\).

Write \(A\) and \(B\) for the even and odd vertices of \(Q_n\). We call
these vertices \emph{roles} or \emph{positions}, and call host vertices
\emph{labels}. Usually the embedding sends \(A\) to \(X\) and \(B\) to
\(Y\); we may interchange the host sides. Distances between cube positions
are Hamming distances.

After labels have been discarded, all size and law normalizations still
use the original side size \(N\), unless another normalization is stated.
For a probability law \(\mu\) on either host side, define its \emph{width}
by
\[
 \operatorname{width}(\mu)=\log\bigl(N\max_x\mu(x)\bigr).
\]
Thus the uniform law on a set of \(M\) labels has width \(\log(N/M)\).
All logarithms are natural.

For a color \(G\) and laws \(\mu\) on \(X\), \(\nu\) on \(Y\), put
\[
 d_G(\mu,\nu)=\sum_{x\in X}\sum_{y\in Y}
                  \mu(x)\nu(y)\mathbf 1_{\{xy\in G\}}.
\]
This is the edge density seen by independent draws from the two laws.
A label in an argument of \(d_G\) denotes its point mass. A label
\emph{hits} another label in \(G\) when the two are adjacent in that
color. We leave integer roundoffs in scales and layer counts implicit
when the accompanying asymptotic inequalities have slack.

The tools in Section~\ref{sec:general-conventions-and-probabilistic-tools}
are used throughout. Sections~\ref{sec:bias-versus-purity-at-power-widths}--\ref{sec:a-small-power-key-grid}
establish the initial discrepancy bound~\eqref{eq:source-2};
Sections~\ref{sec:asymmetric-pure-patches-under-2}--\ref{sec:jump-to-large-bias-only-at-linear-budget}
exclude asymmetric purity, intermediate bias, full-dimensional clusters,
and the linear-budget jump.
Section~\ref{sec:the-two-deep-orientations-some-estimates} records the
resulting deep discrepancy and interaction estimates.
Sections~\ref{sec:patch-tiling-under-9}--\ref{sec:high-modes-of-the-tiling}
extract patches, allocate the cube, and exclude the high modes.
Sections~\ref{sec:low-modes-geometry-and-cell-states}--\ref{sec:late-dynamics-and-its-transfer-problem}
handle the bounded and low modes, ending with the two-endpoint Hall
estimate. Figure~\ref{fig:proof-roadmap} shows the dependencies between
these reductions.
